% ============================================================
%  C: Como Programar -- 6ª Edição | Deitel & Deitel
%  Capítulo 4 -- Controle de Programa em C
%  EXERCÍCIOS DE AUTORREVISÃO (4.1 -- 4.4)
%  Abordagem incremental: Caps. 1 + 2 + 3 + 4
% ============================================================
\documentclass[12pt,a4paper]{article}
\usepackage[utf8]{inputenc}
\usepackage[T1]{fontenc}
\usepackage[brazil]{babel}
\usepackage{geometry}
\geometry{top=2.5cm,bottom=2.5cm,left=3cm,right=3cm}
\usepackage{parskip}\setlength{\parskip}{6pt}\setlength{\parindent}{0pt}
\usepackage{xcolor,tcolorbox,enumitem,listings,fancyhdr,titlesec,hyperref,booktabs,array}
\tcbuselibrary{skins,breakable}

\definecolor{deitelblue}{RGB}{0,70,127}
\definecolor{deitelgray}{RGB}{240,242,245}
\definecolor{deitelgreen}{RGB}{0,110,60}
\definecolor{deitellightblue}{RGB}{220,235,250}
\definecolor{deitelred}{RGB}{160,20,20}
\definecolor{deitellorange}{RGB}{200,90,0}

\newtcolorbox{boaprática}[1][]{colback=deitellightblue,colframe=deitelblue,
  fonttitle=\bfseries\small,title={Boa prática de programação},breakable,#1}
\newtcolorbox{erroco}[1][]{colback=red!5,colframe=deitelred,
  fonttitle=\bfseries\small,title={Erro comum de programação},breakable,#1}
\newtcolorbox{prev}[1][]{colback=yellow!8,colframe=deitellorange,
  fonttitle=\bfseries\small,title={Dica de prevenção de erro},breakable,#1}
\newtcolorbox{respostabox}[1][]{colback=deitelgray,colframe=deitelblue!60,
  fonttitle=\bfseries\small\color{deitelblue},title={Resposta detalhada},breakable,#1}
\newtcolorbox{enunciadoBox}[1][]{colback=white,colframe=deitelblue,
  fonttitle=\bfseries,title={#1},breakable}
\newtcolorbox{saida}[1][]{colback=black!88,colframe=deitelblue,
  fontupper=\ttfamily\small\color{green!70},title={\small\color{white}Saída do programa},breakable,#1}

\setlist[enumerate,1]{label=\textbf{\alph*)},leftmargin=2em}
\setlist[itemize]{leftmargin=2em}

\lstset{
  basicstyle=\ttfamily\small,backgroundcolor=\color{deitelgray},
  frame=single,framerule=0.4pt,rulecolor=\color{deitelblue!40},
  breaklines=true,showstringspaces=false,
  keywordstyle=\color{deitelblue}\bfseries,
  commentstyle=\color{deitelgreen}\itshape,
  stringstyle=\color{deitelred},
  language=C,numbers=left,numberstyle=\tiny\color{gray},
  xleftmargin=18pt,xrightmargin=5pt,stepnumber=1,
}

\pagestyle{fancy}\fancyhf{}
\fancyhead[L]{\small\color{deitelblue}\textbf{C: Como Programar -- 6ª ed. | Deitel \& Deitel}}
\fancyhead[R]{\small\color{deitelblue}Cap. 4 -- Autorrevisão}
\fancyfoot[C]{\small\thepage}
\renewcommand{\headrulewidth}{0.4pt}

\titleformat{\section}[block]{\large\bfseries\color{deitelblue}}{\thesection.}{0.5em}{}[\titlerule]
\titleformat{\subsection}[block]{\normalsize\bfseries\color{deitelblue!80}}{}{}{}

\hypersetup{colorlinks=true,linkcolor=deitelblue}

\begin{document}

\begin{titlepage}
  \centering\vspace*{2cm}
  {\color{deitelblue}\rule{\linewidth}{2pt}}\\[0.4cm]
  {\huge\bfseries\color{deitelblue}C: Como Programar}\\[0.2cm]
  {\large\color{deitelblue}6ª Edição $\cdot$ Paul Deitel \& Harvey Deitel}\\[0.2cm]
  {\color{deitelblue}\rule{\linewidth}{2pt}}\\[1cm]
  {\LARGE\bfseries Capítulo 4}\\[0.3cm]
  {\Large\color{deitelblue!80}Controle de Programa em C}\\[0.8cm]
  \begin{tcolorbox}[colback=deitellightblue,colframe=deitelblue,width=0.85\linewidth]
    \centering{\large\bfseries Exercícios de Autorrevisão (4.1--4.4)}\\[0.2cm]
    {\normalsize Resolvidos passo a passo, abordagem incremental\\
    Capítulos 1 + 2 + 3 + 4}
  \end{tcolorbox}
\end{titlepage}

\section*{Construindo sobre os Capítulos Anteriores}

No Capítulo~4, ampliamos significativamente o controle de programa em C. Construindo sobre os fundamentos
dos Capítulos~1, 2 e 3, incorporamos:

\begin{itemize}[noitemsep]
  \item \textbf{Estrutura \texttt{for}} -- ideal para repetição controlada por contador
  \item \textbf{Estrutura \texttt{do...while}} -- testa a condição \textit{depois} do corpo
  \item \textbf{Estrutura \texttt{switch}} -- seleção múltipla baseada em valores constantes
  \item \textbf{\texttt{break} e \texttt{continue}} -- modificadores do fluxo de loops
  \item \textbf{Operadores lógicos} \texttt{\&\&} (E), \texttt{||} (OU), \texttt{!} (NÃO)
  \item \textbf{Diferença entre \texttt{=} e \texttt{==}} -- atribuição vs.\ comparação
  \item \textbf{Função \texttt{pow}} da biblioteca matemática \texttt{<math.h>}
  \item \textbf{Formatação avançada de \texttt{printf}} -- largura de campo e alinhamento
\end{itemize}

\begin{boaprática}
  Os recursos do Capítulo~4 dão ao seu programa muito mais expressividade. \texttt{for}
  é a forma idiomática de loops contados; \texttt{switch} substitui cadeias de
  \texttt{if...else} quando comparamos uma variável com valores constantes; e os
  operadores lógicos permitem combinar condições naturalmente.
\end{boaprática}

\tableofcontents\newpage

% ============================================================
\section{Exercício 4.1 -- Preenchimento de Lacunas}
% ============================================================

\begin{enunciadoBox}[Enunciado 4.1]
Preencha os espaços nas seis sentenças sobre repetição controlada por contador,
repetição controlada por sentinela, comandos de controle e estrutura \texttt{switch}.
\end{enunciadoBox}

\subsection*{Item (a) -- Repetição controlada por contador}

\begin{respostabox}
\textbf{Resposta:} \textcolor{deitelblue}{\textbf{definida}}

A repetição controlada por contador é também conhecida como \textit{repetição definida},
porque o número de iterações é \textbf{conhecido com antecedência}, antes que o laço
comece a ser executado. Por exemplo, ``execute esta tarefa exatamente 10 vezes''.
A estrutura \texttt{for} foi projetada especificamente para esse padrão.
\end{respostabox}

\subsection*{Item (b) -- Repetição controlada por sentinela}

\begin{respostabox}
\textbf{Resposta:} \textcolor{deitelblue}{\textbf{indefinida}}

A repetição controlada por sentinela é também conhecida como \textit{repetição indefinida},
porque \textbf{não sabemos com antecedência} quantas vezes o laço executará. O usuário
controla isso digitando o valor de sentinela quando desejar terminar.
\end{respostabox}

\subsection*{Item (c) -- O elemento contador da repetição definida}

\begin{respostabox}
\textbf{Resposta:} \textcolor{deitelblue}{\textbf{variável de controle}} (ou \textbf{contador})

A \textit{variável de controle} é a variável que conta o número de iterações executadas.
Em uma estrutura \texttt{for}, ela aparece três vezes no cabeçalho:
\begin{lstlisting}
for ( i = 1; i <= 10; ++i ) {
    /* corpo do for */
}
\end{lstlisting}
\begin{itemize}[noitemsep]
  \item \texttt{i = 1} -- valor \textit{inicial} da variável de controle
  \item \texttt{i <= 10} -- condição de continuação testada \textit{antes} de cada iteração
  \item \texttt{++i} -- incremento da variável de controle, após cada iteração
\end{itemize}
\end{respostabox}

\subsection*{Item (d) -- Comando para a próxima iteração}

\begin{respostabox}
\textbf{Resposta:} \textcolor{deitelblue}{\textbf{\texttt{continue}}}

O comando \texttt{continue} \textit{pula} o restante do corpo do laço e prossegue
imediatamente com a próxima iteração. Em \texttt{while} e \texttt{do...while}, isso
significa retornar ao teste da condição; em \texttt{for}, executa primeiro a expressão
de incremento.

\begin{lstlisting}
/* Imprime 1 a 10, mas pula 5 */
for ( i = 1; i <= 10; ++i ) {
    if ( i == 5 ) {
        continue;  /* pula esta iteracao */
    }
    printf( "%d ", i );
}
/* Saida: 1 2 3 4 6 7 8 9 10 */
\end{lstlisting}
\end{respostabox}

\subsection*{Item (e) -- Comando para sair imediatamente da estrutura}

\begin{respostabox}
\textbf{Resposta:} \textcolor{deitelblue}{\textbf{\texttt{break}}}

O comando \texttt{break} causa a saída \textbf{imediata} de uma estrutura de repetição
(\texttt{for}, \texttt{while}, \texttt{do...while}) ou de uma estrutura \texttt{switch}.
A execução continua na primeira instrução \textit{após} a estrutura.

\begin{lstlisting}
/* Encontra o primeiro multiplo de 7 acima de 50 */
for ( i = 50; i <= 100; ++i ) {
    if ( i % 7 == 0 ) {
        printf( "Encontrado: %d\n", i );
        break;  /* sai do for imediatamente */
    }
}
\end{lstlisting}
\end{respostabox}

\subsection*{Item (f) -- Estrutura para testar valores inteiros constantes}

\begin{respostabox}
\textbf{Resposta:} \textcolor{deitelblue}{\textbf{estrutura de seleção múltipla \texttt{switch}}}

A estrutura \texttt{switch} (Seção~4.7) testa uma variável (ou expressão) contra
valores inteiros constantes -- chamados \textit{rótulos \texttt{case}}. Para cada
rótulo correspondente, executa as instruções associadas até encontrar um \texttt{break}.
O caso \texttt{default} (opcional) captura todos os valores não casados.

\begin{lstlisting}
switch ( opcao ) {
    case 1:
        printf( "Opcao 1 escolhida\n" );
        break;
    case 2:
        printf( "Opcao 2 escolhida\n" );
        break;
    default:
        printf( "Opcao invalida\n" );
        break;
}
\end{lstlisting}
\end{respostabox}

\newpage

% ============================================================
\section{Exercício 4.2 -- Verdadeiro ou Falso}
% ============================================================

\subsection*{Item (a) -- O caso default é obrigatório no switch}

\begin{tcolorbox}[colback=red!5,colframe=deitelred,title={\bfseries FALSO}]
O caso \texttt{default} é \textbf{opcional}. Se não for necessário tomar nenhuma
ação para valores não casados pelos rótulos \texttt{case}, o \texttt{default} pode
ser omitido. Sem \texttt{default}, valores não correspondentes simplesmente não
disparam nenhuma ação.

\begin{boaprática}
  Mesmo quando aparentemente desnecessário, fornecer um \texttt{default} é uma
  excelente prática defensiva. Ele captura valores inesperados (devido a bugs
  ou entradas inválidas) que poderiam passar silenciosamente despercebidos.
\end{boaprática}
\end{tcolorbox}

\subsection*{Item (b) -- O \texttt{break} é obrigatório no caso default}

\begin{tcolorbox}[colback=red!5,colframe=deitelred,title={\bfseries FALSO}]
O \texttt{break} é \textbf{usado para sair da estrutura \texttt{switch}}, evitando
o efeito conhecido como \textit{fall-through} (queda em cascata para o próximo case).
Quando o \texttt{default} é o \textit{último} rótulo do \texttt{switch}, o \texttt{break}
é desnecessário, pois não há mais casos para os quais ``cair''.

Ainda assim, é uma boa prática colocar \texttt{break} no \texttt{default}, pois se
um case adicional for adicionado depois, no futuro, o \texttt{break} já estará lá
prevenindo bugs.
\end{tcolorbox}

\subsection*{Item (c) -- A expressão \texttt{(x > y \&\& a < b)}}

\begin{tcolorbox}[colback=red!5,colframe=deitelred,title={\bfseries FALSO}]
O operador \texttt{\&\&} (E lógico) exige que \textbf{ambas} as expressões relacionais
sejam verdadeiras para que o resultado total seja verdadeiro. A enunciação ``\textit{ou}''
do enunciado descreve o operador \texttt{||} (OU lógico), não o \texttt{\&\&}.

\textbf{Tabela-verdade do \texttt{\&\&}:}
\begin{center}
\renewcommand{\arraystretch}{1.3}
\begin{tabular}{c c c}
  \toprule
  \textbf{x > y} & \textbf{a < b} & \textbf{x > y \&\& a < b} \\ \midrule
  falso & falso & falso \\
  falso & verdadeiro & falso \\
  verdadeiro & falso & falso \\
  verdadeiro & verdadeiro & \textbf{verdadeiro} \\
  \bottomrule
\end{tabular}
\end{center}
\end{tcolorbox}

\subsection*{Item (d) -- A expressão com \texttt{||}}

\begin{tcolorbox}[colback=green!5,colframe=deitelgreen,title={\bfseries VERDADEIRO}]
O operador \texttt{||} (OU lógico) é verdadeiro se \textbf{pelo menos um} dos
operandos for verdadeiro. Em particular, é verdadeiro também quando \textit{ambos}
são verdadeiros. Apenas quando \textit{ambos} são falsos é que o \texttt{||}
produz falso.
\end{tcolorbox}

\newpage

% ============================================================
\section{Exercício 4.3 -- Instruções com novos recursos}
% ============================================================

\subsection*{Item (a) -- Soma dos ímpares de 1 a 99 com \texttt{for}}

\begin{respostabox}
\begin{lstlisting}
soma = 0;
for ( contador = 1; contador <= 99; contador += 2 ) {
    soma += contador;
}
\end{lstlisting}
\textbf{Explicação:} Inicializamos \texttt{contador} em 1 (primeiro ímpar) e
incrementamos de 2 em 2 (\texttt{contador += 2}), garantindo que percorremos
apenas ímpares: 1, 3, 5, \ldots, 99. \textbf{Resultado:} $1 + 3 + 5 + \ldots + 99 = 2500$.
\end{respostabox}

\subsection*{Item (b) -- Imprimir 333,546372 com 5 precisões diferentes}

\begin{respostabox}
\begin{lstlisting}
printf( "%-15.1f\n", 333.546372 );  /* imprime: 333.5            */
printf( "%-15.2f\n", 333.546372 );  /* imprime: 333.55           */
printf( "%-15.3f\n", 333.546372 );  /* imprime: 333.546          */
printf( "%-15.4f\n", 333.546372 );  /* imprime: 333.5464         */
printf( "%-15.5f\n", 333.546372 );  /* imprime: 333.54637        */
\end{lstlisting}
\textbf{Explicação dos especificadores:}
\begin{itemize}[noitemsep]
  \item \texttt{\%} -- início do especificador
  \item \texttt{-} -- alinha à \textbf{esquerda} no campo
  \item \texttt{15} -- largura do campo (15 posições)
  \item \texttt{.N} -- precisão (N casas decimais)
  \item \texttt{f} -- tipo \texttt{float} ou \texttt{double}
\end{itemize}

\textbf{Os 5 valores impressos:} \texttt{333.5}, \texttt{333.55}, \texttt{333.546},
\texttt{333.5464}, \texttt{333.54637}. Note o arredondamento, não truncamento.
\end{respostabox}

\subsection*{Item (c) -- Calcular $2{,}5^3$ usando \texttt{pow}}

\begin{respostabox}
\begin{lstlisting}
#include <math.h>  /* necessario para usar pow */
/* ... */
printf( "%10.2f\n", pow( 2.5, 3 ) );  /* imprime: "    15.63" */
\end{lstlisting}

\textbf{Valor impresso:} \texttt{15.63} (com 6 espaços antes, totalizando 10 posições).

\textbf{Observações:}
\begin{itemize}[noitemsep]
  \item \texttt{pow(base, expoente)} retorna $\text{base}^{\text{expoente}}$ como \texttt{double}
  \item $2{,}5^3 = 2{,}5 \times 2{,}5 \times 2{,}5 = 15{,}625$ \(\to\) arredondado a 2 casas: \texttt{15.63}
  \item Sem o \texttt{-}, o valor é alinhado à \textbf{direita} no campo de 10 caracteres
\end{itemize}

\begin{boaprática}
  Para usar funções da biblioteca matemática (\texttt{pow}, \texttt{sqrt}, \texttt{sin},
  etc.), inclua \texttt{<math.h>}. Em sistemas Linux, ao compilar, pode ser necessário
  ligar a biblioteca matemática com \texttt{-lm}: \texttt{gcc programa.c -o programa -lm}.
\end{boaprática}
\end{respostabox}

\subsection*{Item (d) -- Imprimir 1 a 20 com \texttt{while}, 5 por linha}

\begin{respostabox}
\begin{lstlisting}
x = 1;
while ( x <= 20 ) {
    printf( "%d", x );
    if ( x % 5 == 0 ) {
        printf( "\n" );  /* nova linha a cada 5 numeros */
    }
    else {
        printf( "\t" );  /* tabulacao entre numeros     */
    }
    ++x;
}
\end{lstlisting}

\textbf{Saída esperada:}
\begin{saida}
1\textbackslash{}t2\textbackslash{}t3\textbackslash{}t4\textbackslash{}t5\\
6\textbackslash{}t7\textbackslash{}t8\textbackslash{}t9\textbackslash{}t10\\
11\textbackslash{}t12\textbackslash{}t13\textbackslash{}t14\textbackslash{}t15\\
16\textbackslash{}t17\textbackslash{}t18\textbackslash{}t19\textbackslash{}t20
\end{saida}

\textbf{Lógica do truque:} \texttt{x \% 5 == 0} é verdadeiro apenas para múltiplos
de 5 (5, 10, 15, 20) -- exatamente onde queremos quebrar a linha.
\end{respostabox}

\subsection*{Item (e) -- O mesmo, com \texttt{for}}

\begin{respostabox}
\begin{lstlisting}
for ( x = 1; x <= 20; ++x ) {
    printf( "%d", x );
    if ( x % 5 == 0 ) {
        printf( "\n" );
    }
    else {
        printf( "\t" );
    }
}
\end{lstlisting}

\textbf{Comparação \texttt{while} vs.\ \texttt{for}:} A versão \texttt{for} é mais
\textit{idiomática} para repetição controlada por contador. As três partes do controle
(inicialização, condição, incremento) ficam concentradas no cabeçalho do \texttt{for},
tornando o código mais conciso e legível.
\end{respostabox}

\newpage

% ============================================================
\section{Exercício 4.4 -- Identificar e Corrigir Erros}
% ============================================================

\subsection*{Item (a) -- Loop infinito por \texttt{;} extra}

\begin{respostabox}
\textbf{Código com erro:}
\begin{lstlisting}
x = 1;
while ( x <= 10 );
    x++;
}
\end{lstlisting}

\textbf{Erros identificados:}
\begin{enumerate}[noitemsep]
  \item O ponto e vírgula após \texttt{while ( x <= 10 )} cria um corpo de laço \textit{vazio};
  como \texttt{x} nunca é incrementado, gera-se um \textbf{loop infinito}.
  \item Falta a chave de abertura \texttt{\{} (há \texttt{\}} no final, mas sem \texttt{\{}).
\end{enumerate}

\textbf{Correção:}
\begin{lstlisting}
x = 1;
while ( x <= 10 ) {  /* remove o ; e abre chave aqui */
    x++;
}
\end{lstlisting}
\end{respostabox}

\subsection*{Item (b) -- Ponto flutuante como variável de controle}

\begin{respostabox}
\textbf{Código com erro:}
\begin{lstlisting}
for ( y = .1; y != 1.0; y += .1 )
    printf( "%f\n", y );
\end{lstlisting}

\textbf{Erro:} Usar um número de \textbf{ponto flutuante} como variável de controle de
um \texttt{for} é \textit{perigosíssimo}. Devido à imprecisão da representação binária
de números fracionários (você não consegue representar 0,1 exatamente em binário, da
mesma forma que não consegue representar 1/3 exatamente em decimal), a soma cumulativa
de incrementos \texttt{0.1} \textbf{não chega} exatamente a \texttt{1.0}, e a condição
\texttt{y != 1.0} nunca se torna verdadeira. \textbf{Loop infinito}!

\textbf{Correção:} use um inteiro como variável de controle e faça a divisão
para obter o valor fracionário desejado.
\begin{lstlisting}
for ( y = 1; y != 10; y++ )
    printf( "%f\n", (float) y / 10 );
\end{lstlisting}

\begin{prev}
  \textbf{Regra de ouro:} \textit{nunca} compare valores de ponto flutuante por
  igualdade exata (\texttt{==} ou \texttt{!=}). Se for absolutamente necessário,
  teste se a diferença é menor que uma pequena tolerância:
  \texttt{fabs(x - y) < 1e-9}.
\end{prev}
\end{respostabox}

\subsection*{Item (c) -- Falta \texttt{break} no \texttt{switch}}

\begin{respostabox}
\textbf{Código com erro:}
\begin{lstlisting}
switch ( n ) {
    case 1:
        printf( "O numero e 1\n" );
    case 2:
        printf( "O numero e 2\n" );
        break;
    default:
        printf( "O numero nao e nem 1 nem 2\n" );
        break;
}
\end{lstlisting}

\textbf{Erro:} Falta o \texttt{break} no final do \texttt{case 1}. Quando \texttt{n == 1},
o programa imprime ``O numero e 1'' e \textbf{continua executando} as instruções do
\texttt{case 2} (\textit{fall-through}), imprimindo também ``O numero e 2'' --
comportamento incorreto.

\textbf{Correção:}
\begin{lstlisting}
switch ( n ) {
    case 1:
        printf( "O numero e 1\n" );
        break;  /* adiciona break para evitar fall-through */
    case 2:
        printf( "O numero e 2\n" );
        break;
    default:
        printf( "O numero nao e nem 1 nem 2\n" );
        break;
}
\end{lstlisting}

\textit{Nota:} Em raras situações, o \textit{fall-through} é \textbf{intencional}.
Por exemplo, na contagem de notas (Figura~4.7 do livro), \texttt{case 'A'} cai em
\texttt{case 'a'} para tratar maiúscula e minúscula igualmente.
\end{respostabox}

\subsection*{Item (d) -- Operador relacional incorreto}

\begin{respostabox}
\textbf{Código com erro:}
\begin{lstlisting}
n = 1;
while ( n < 10 )
    printf( "%d ", n++ );
\end{lstlisting}

\textbf{Erro:} A condição \texttt{n < 10} é \textit{estritamente menor que}, então o
laço termina quando \texttt{n} chega a 10 -- imprimindo apenas de 1 a 9, não 1 a 10.

\textbf{Correção:} use \texttt{<=} em vez de \texttt{<}.
\begin{lstlisting}
n = 1;
while ( n <= 10 )
    printf( "%d ", n++ );
\end{lstlisting}

\begin{erroco}
  Erros de \textbf{off-by-one} (``erro de uma posição'') -- usar \texttt{<} em vez
  de \texttt{<=} ou vice-versa -- são entre os mais comuns e perigosos. Eles
  podem fazer um laço executar uma vez a menos (ou a mais) do que deveria,
  com resultados sutis difíceis de detectar.
\end{erroco}
\end{respostabox}

\newpage

% ============================================================
\section*{Gabarito Rápido -- Capítulo 4: Autorrevisão}

\begin{tcolorbox}[colback=deitellightblue,colframe=deitelblue,title={\bfseries\large Respostas Resumidas}]

\textbf{4.1:}
\begin{itemize}[noitemsep]
  \item[a)] definida \quad b) indefinida \quad c) variável de controle (contador)
  \item[d)] \texttt{continue} \quad e) \texttt{break} \quad f) \texttt{switch}
\end{itemize}

\medskip\textbf{4.2 -- Verdadeiro/Falso:}
\begin{itemize}[noitemsep]
  \item[a)] \textcolor{deitelred}{\textbf{Falso}} -- \texttt{default} é opcional
  \item[b)] \textcolor{deitelred}{\textbf{Falso}} -- \texttt{break} pode ser dispensado se \texttt{default} for o último
  \item[c)] \textcolor{deitelred}{\textbf{Falso}} -- \texttt{\&\&} exige ambas verdadeiras
  \item[d)] \textcolor{deitelgreen}{\textbf{Verdadeiro}} -- \texttt{||} é verdadeiro se um ou ambos forem verdadeiros
\end{itemize}

\medskip\textbf{4.3:} Soluções com \texttt{for}, \texttt{while}, \texttt{pow} e formatação \texttt{\%-15.Nf}.

\medskip\textbf{4.4:}
\begin{itemize}[noitemsep]
  \item[a)] \texttt{;} extra após \texttt{while} cria loop infinito
  \item[b)] Não use \texttt{float} como controle de \texttt{for}
  \item[c)] Falta \texttt{break} no \texttt{case 1}
  \item[d)] Use \texttt{<=} em vez de \texttt{<}
\end{itemize}

\end{tcolorbox}

\end{document}
